\chapter{Complete Coefficients on the Pro-etale Site}
\label{chap:proetale-coefficients}

Fix a prime \(\ell\) invertible on \(X\).  The basic coefficient rings will
be \(\Lambda=\Zl\) and, more generally, the ring of integers
\(\Lambda=\mathcal O_E\) in a finite extension \(E/\Ql\).  Write
\(\mathfrak m=(\varpi)\) and \(\Lambda_n=\Lambda/\mathfrak m^n\).  These are
complete discrete valuation rings, so the ideal of definition is regular.
Regularity permits the unbounded system comparison in
\cite[Lemma~3.5.7(5)]{BhattScholzeProEtale}.  Repleteness gives
\(\R\varprojlim\) cohomological dimension at most one, while surjective
reduction towers satisfy the Mittag--Leffler condition and kill the relevant
\(\varprojlim^1\).  No exactness claim is made for arbitrary inverse systems.

\section{Derived completion}

Let \(\underline\Lambda\) denote the continuous coefficient sheaf of
\cref{prop:proetale-zl-sheaf}, with the identical construction for
\(\mathcal O_E\).

\begin{definition}[Derived complete complex]
\label{def:derived-complete-proetale}
For \(K\in\D(X_{\proet},\underline\Lambda)\), put
\[
  K_n=K\mathbin{\mathop\otimes\limits^{\mathbf L}_{\Lambda}}\Lambda_n.
\]
The complex \(K\) is \emph{derived \(\mathfrak m\)-complete} if the canonical
map
\begin{equation}
\label{eq:derived-completion-map}
 K\longrightarrow
 \Rlim_n K_n
\end{equation}
is an equivalence.  Its derived completion is
\[
  \widehat K:=\Rlim_n
  \bigl(K\mathbin{\mathop\otimes\limits^{\mathbf L}_{\Lambda}}\Lambda_n\bigr).
\]
Write \(\Dcomp(X,\Lambda)\) for the full subcategory of complete objects.
\end{definition}

\begin{proposition}[Completion is a reflector]
\label{prop:completion-reflector}
The inclusion
\[
  \Dcomp(X,\Lambda)\hookrightarrow
  \D(X_{\proet},\underline\Lambda)
\]
admits \(K\mapsto\widehat K\) as a left adjoint.  For complete \(K,L\), define
\[
 K\widehat\otimes^{\mathbf L}_{\Lambda}L
   :=\widehat{K\otimes^{\mathbf L}_{\Lambda}L},
 \qquad
 \R\underline{\operatorname{Hom}}_{\Lambda}(K,L)
\]
with the internal Hom taken in complete objects.  These operations make
\(\Dcomp(X,\Lambda)\) a closed symmetric monoidal stable category.
\end{proposition}

\begin{proof}
For a finitely generated ideal in a replete topos, the abstract completion
reflector is \cite[Lemma~3.4.9]{BhattScholzeProEtale}.  In the noetherian
constant-ring case, \cite[Proposition~3.5.1]{BhattScholzeProEtale} identifies
that reflector with \(K\mapsto\Rlim_n(K\otimes^{\mathbf L}\Lambda_n)\) and
proves the adjunction.  The completed coefficient sheaf and its complete
module category are \cite[Definition~3.5.2]{BhattScholzeProEtale}.
Completed tensor is \cite[Definition~3.4.10]{BhattScholzeProEtale}, while
closedness and tensor--Hom adjunction are
\cite[Lemma~3.4.11]{BhattScholzeProEtale}.  Finally,
\cite[Lemma~3.5.7]{BhattScholzeProEtale} supplies the equivalence with
compatible systems used in the next section; its unbounded part uses that
\((\varpi)\) is a regular ideal.
\end{proof}

\section{Complete objects and compatible torsion systems}

\begin{definition}[Compatible system]
A compatible \(\Lambda_n\)-system is a sequence
\[
 K_n\in\D(X_{\proet},\Lambda_n)
\]
equipped with coherent equivalences
\[
 K_{n+1}\otimes^{\mathbf L}_{\Lambda_{n+1}}\Lambda_n
   \simeq K_n.
\]
Morphisms must respect all transition equivalences and their higher
coherences.
\end{definition}

\begin{theorem}[Completion--system equivalence]
\label{thm:proetale-complete-system-equivalence}
Reduction modulo \(\mathfrak m^n\) induces an equivalence of stable
categories
\begin{equation}
\label{eq:completion-system-equivalence}
 \Dcomp(X,\Lambda)
   \xrightarrow{\sim}
 \varprojlim_n \D(X_{\proet},\Lambda_n).
\end{equation}
Its inverse sends \((K_n)_n\) to \(\Rlim_n K_n\).  The statement remains
valid in the unbounded derived category for \(\Zl\) and \(\mathcal O_E\).
\end{theorem}

\begin{proof}
For complete \(K\), \cref{eq:derived-completion-map} identifies \(K\) with
the inverse limit of its reductions, and the reduction identities supply a
compatible system.  Conversely, let \(K=\Rlim K_n\).  Exactness of products
and the cohomological-dimension-one property of \(\Rlim\) allow derived
tensor with \(\Lambda_m\) to pass through the tower; compatibility makes the
result \(K_m\).  Thus \(K\) is complete and the two constructions are inverse.
The regularity of \((\varpi)\) supplies the unbounded assertion.  This is
\cite[Lemma~3.5.7]{BhattScholzeProEtale}; a coefficient-system formulation is
\cite[Proposition~5.1]{HRSConstructible}.
\end{proof}

\begin{proetalebox}
Traditional adic formalisms begin with a tower
\((K_n)_n\) and then quotient by essentially zero systems.  On the pro-etale
site the tower has a genuine inverse limit, and
\cref{thm:proetale-complete-system-equivalence} says that no information is
lost by replacing the tower with one complete object.  This is the first
major simplification used later.
\end{proetalebox}

\section{Lisse and constructible complexes}

\begin{definition}[Lisse and constructible]
\label{def:proetale-lisse-constructible}
A complete complex \(K\) is \emph{lisse} if it is dualizable for
\(\widehat\otimes^{\mathbf L}_{\Lambda}\).  It is \emph{constructible} if
there is a finite stratification
\[
  X=\coprod_{a=1}^r X_a
\]
by locally closed constructible subschemes such that \(K|_{X_a}\) is lisse
for every \(a\).  We denote the corresponding categories by
\(\D_{\lis}(X,\Lambda)\) and \(\Dcons(X,\Lambda)\).
\end{definition}

For noetherian \(X\), this agrees with the Bhatt--Scholze definition:
\(K\) is complete and \(K\otimes^{\mathbf L}_{\Lambda}\Lambda_1\) has bounded
constructible cohomology.  The equivalence is
\cite[Definition~6.5.1, Lemma~6.5.3, and Proposition~6.6.11]
{BhattScholzeProEtale}.

\begin{theorem}[Local perfect-complex criterion]
\label{thm:proetale-constructible-characterization}
Let \(X\) be locally noetherian.
\begin{enumerate}
\item A complete complex is lisse if and only if pro-etale locally it is the
      constant sheaf attached to a perfect \(\Lambda\)-complex.
\item On a w-contractible affine \(U\), the section functor identifies lisse
      objects with perfect complexes over
      \(\Gamma(U,\underline\Lambda)\).
\item Lisse and constructible complexes satisfy pro-etale hyperdescent.
\item If \(X\) is also qcqs, constructible complexes are bounded.  Completed
      tensor products and internal Homs preserve constructibility.
\end{enumerate}
\end{theorem}

\begin{proof}
On a w-contractible affine the section equivalence is
\cite[Lemma~4.1]{HRSConstructible}, and hyperdescent is
\cite[Corollary~4.7]{HRSConstructible}.  The converse from dualizable to
pro-etale-locally perfect-constant requires that \(X\) have locally finitely
many irreducible components; this is
\cite[Theorem~4.13]{HRSConstructible}.  A locally noetherian scheme satisfies
that hypothesis.  Under qcqs, finitely many affine pieces and their finite
stratifications give a uniform cohomological bound.  Tensor and internal Hom
reduce on each stratum to perfect-complex algebra and then glue through
localization triangles.  Compare
\cite[Lemmas~6.5.3 and 6.5.5]{BhattScholzeProEtale}.
\end{proof}

\begin{example}[The point]
\label{ex:constructible-point}
For a separably closed field \(k\),
\[
  \Dcons(\Spec k,\Lambda)\simeq\operatorname{Perf}(\Lambda).
\]
For a general field, the right-hand side is replaced by perfect
\(\Lambda\)-complexes with continuous action of its absolute Galois group.
The first assertion is \cite[Example~6.5.2]{BhattScholzeProEtale}.
\end{example}

\section{The standard t-structure and a finite-level trap}

\begin{theorem}[Adic t-structure]
\label{thm:proetale-adic-tstructure}
Let \(\Lambda\) be t-admissible: its underlying ring \(\Lambda_*\) is regular
coherent, and
\[
 \Lambda_*\longrightarrow\Gamma(S,\Lambda)
\]
is flat for every extremally disconnected profinite set \(S\).
\begin{enumerate}
\item If \(X\) has Zariski-locally finitely many irreducible components, the
      ambient standard t-structure restricts to
      \(\D_{\lis}(X,\Lambda)\).
\item If every constructible subset of \(X\) has Zariski-locally finitely
      many irreducible components, it restricts to
      \(\Dcons(X,\Lambda)\).
\end{enumerate}
Both geometric conditions hold when \(X\) is locally noetherian.
The topological rings \(\Zl\) and \(\mathcal O_E\) are t-admissible.  Hence,
under the second geometric condition,
\[
  \cH^i(K)\in\Dcons(X,\Lambda)
  \qquad(K\in\Dcons(X,\Lambda)).
\]
\end{theorem}

\begin{proof}
By \cite[Theorem~4.13]{HRSConstructible}, lisse objects are locally
perfect-constant under the first geometric hypothesis.  Regular coherence of
\(\Lambda_*\) keeps truncations of the perfect coefficient complex perfect.
The second part of t-admissibility is equally necessary: flatness of
\(\Lambda_*\to\Gamma(S,\Lambda)\) makes truncation commute with evaluation on
w-contractible affines, whose component spaces are extremally disconnected.
The constructible assertion follows by applying this calculation on a
stratification; the stated condition on every constructible subset ensures
that the calculation remains available on all strata.  This is
\cite[Theorem~6.2]{HRSConstructible}.  Complete discrete valuation rings and
their fraction fields are covered by
\cite[Lemma~6.10 and Corollary~6.11]{HRSConstructible}.
\end{proof}

\begin{warning}[Do not take the inverse limit of finite-level hearts]
For \(n>1\), the ring \(\mathbf Z/\ell^n\mathbf Z\) is not regular.  A truncation of a
perfect complex over it need not be perfect; for example the residue module
\(\Fl\) has infinite projective dimension.  Hence
\(\operatorname{Perf}(\mathbf Z/\ell^n\mathbf Z)\) is not preserved by the standard truncation functors.
The integral t-structure must be constructed on the complete
\(\Zl\)-category and cannot be defined naively as the inverse limit of
finite-level hearts.  This is the warning in
\cite[Section~6]{HRSConstructible}.
\end{warning}

\section{The coefficient tower}

Choose a finite extension \(E/\Ql\), with ring of integers \(\mathcal O_E\).

\begin{theorem}[Integral, rational, and algebraically closed coefficients]
\label{thm:proetale-coefficient-tower}
The constructible categories are connected by the following operations.
\begin{enumerate}
\item Derived completed scalar extension gives
\[
  -\widehat\otimes^{\mathbf L}_{\Zl}\mathcal O_E:
  \Dcons(X,\Zl)\longrightarrow\Dcons(X,\mathcal O_E).
\]
This is a change of coefficients, not an equivalence.
\item Inverting \(\varpi\) induces an equivalence
\begin{equation}
\label{eq:oe-to-e}
  \Dcons(X,\mathcal O_E)[1/\ell]
     \xrightarrow{\sim}\Dcons(X,E).
\end{equation}
\item The algebraic-closure coefficient category is the filtered colimit
\begin{equation}
\label{eq:qbarl-colimit}
  \Dcons(X,\Qbarl)
    :=\varinjlim_{E/\Ql\,, [E:\Ql]<\infty}\Dcons(X,E),
\end{equation}
with transition functors given by scalar extension.
\end{enumerate}
Every constructible \(E\)-complex admits an \(\mathcal O_E\)-model after
possibly enlarging \(E\); such a lattice is not canonical.
\end{theorem}

\begin{proof}
Completed scalar extension may be checked after reduction modulo powers of
the maximal ideals, where it preserves perfectness on strata.  For the second
assertion, an integral object killed after inverting \(\ell\) is precisely
\(\ell\)-power torsion, and every rational constructible object admits a
lattice locally on a finite stratification.  Gluing by localization triangles
gives essential surjectivity; clearing denominators gives full faithfulness.
These are \cite[Proposition~6.8.4, Lemma~6.8.13, and
Proposition~6.8.14]{BhattScholzeProEtale}.  The colimit description is
\cite[Lemma~7.9]{HRSConstructible}.
\end{proof}

\begin{remark}[What \(\Qbarl\) means]
Equation~\eqref{eq:qbarl-colimit} is deliberately categorical.  It avoids
treating the non-discrete, non-locally compact field \(\Qbarl\) as though it
were another complete discrete valuation ring.  Every individual
constructible object and every morphism is defined over some finite
extension \(E/\Ql\).
\end{remark}

\begin{theorem}[Comparison with the classical adic category]
\label{thm:proetale-classical-adic-comparison}
For a noetherian scheme on which \(\ell\) is invertible, the constructible
pro-etale categories with \(\mathcal O_E\)- and \(E\)-coefficients agree with
the classical bounded constructible adic derived categories.  Under this
equivalence, reduction modulo \(\mathfrak m^n\), tensor products, pullback,
and cohomology agree.
\end{theorem}

\begin{proof}
Apply \cref{thm:proetale-complete-system-equivalence} to replace a complete
object by its compatible torsion tower.  The comparison of discrete
coefficients in \cref{prop:proetale-etale-comparison} identifies every level
with the classical etale category.  The adic normalization condition is
exactly the derived reduction compatibility of the tower.  The complete
comparison, including rational localization, is
\cite[Theorem~7.7]{HRSConstructible}.
\end{proof}

\begin{formalization}
The reusable API should distinguish:
\[
\begin{array}{c}
\text{all pro-etale complexes}\\
\downarrow;\text{solidification}\\
\text{solid ambient complexes}\\
\supset \text{ind-constructible complexes}\\
\supset \text{constructible complexes}\\
\supset \text{lisse complexes}.
\end{array}
\]
Derived completion remains a separate reflector on the integral coefficient
side.  The system comparison is an equivalence of stable categories, not an
equality of objects.  Solidification, completion, scalar extension, rational
coefficients, and choice of a lattice must be separate constructions; no
theorem may depend on a definitional choice of lattice.
\end{formalization}

\begin{auditbox}
This chapter constructs the coefficient categories and proves their
comparison with traditional adic systems.  It does not prove that geometric
direct images preserve constructibility.  The next chapter first constructs
the operations on the solid ambient category and then proves their
constructible restriction.  The latter still depends on Grothendieck--Gabber
finiteness and constructible duality.
\end{auditbox}

